\documentclass[12pt]{article}
\usepackage{amsmath, amssymb, amsthm}
\usepackage{mathtools}
\usepackage{tensor}
\usepackage{geometry}
\usepackage{booktabs}
\usepackage{array}
\usepackage{xcolor}
\usepackage{hyperref}
\usepackage{fancyhdr}
\usepackage{slashed}
\geometry{margin=1.1in, headheight=15pt}

\definecolor{cadblue}{RGB}{30,90,200}
\definecolor{verifygreen}{RGB}{20,140,60}

\newcommand{\cdbexpr}[1]{\textcolor{cadblue}{$#1$}}
\newcommand{\verified}[1]{\textcolor{verifygreen}{\checkmark\; #1}}
\newcommand{\half}{\tfrac{1}{2}}
\newcommand{\dal}{\dot{\alpha}}
\newcommand{\dbe}{\dot{\beta}}
\newcommand{\dga}{\dot{\gamma}}
\newcommand{\dde}{\dot{\delta}}
\newcommand{\sigmabar}{\bar{\sigma}}

\pagestyle{fancy}
\fancyhf{}
\lhead{Srednicki QFT --- Chapter 60}
\rhead{Spinor Helicity for Spinor QED}
\cfoot{\thepage}

\title{\textbf{Srednicki QFT: Chapter 60}\\[6pt]
       \large Spinor Helicity for Spinor Electrodynamics\\[4pt]
       \normalsize Cadabra2 expressions and numerical verification}
\author{Generated by \texttt{ch60\_export\_latex.py}}
\date{}

\begin{document}
\maketitle
\tableofcontents
\newpage

%% ============================================================
\section{§60.A\quad Twistor Notation}
%% ============================================================

\subsection{The four twistors}

For each massless 4-momentum $p^\mu$ (with $p^2=0$), Srednicki defines
four \textbf{twistors} (massless bispinors, eq.~60.1):
\begin{align}
  |p] &\equiv u_-(p) = v_+(p)
  &&\text{positive helicity ket (square bracket)}
  \tag{60.1a}\\
  |p\rangle &\equiv u_+(p) = v_-(p)
  &&\text{negative helicity ket (angle bracket)}
  \tag{60.1b}\\
  [p| &\equiv \bar u_+(p) = \bar v_-(p)
  &&\text{positive helicity bra (square bracket)}
  \tag{60.1c}\\
  \langle p| &\equiv \bar u_-(p) = \bar v_+(p)
  &&\text{negative helicity bra (angle bracket)}
  \tag{60.1d}
\end{align}

\textbf{Physical interpretation.}
For a massless particle, the Dirac equation $\slashed{p}\,\psi=0$ has the same
mathematical solutions regardless of whether $\psi$ is interpreted as a particle
($u$) or antiparticle ($v$).  The identification $u_-(p)=v_+(p)$ encodes the
\emph{crossing symmetry} of the $S$-matrix: flipping an incoming particle to
an outgoing antiparticle corresponds to $u\leftrightarrow v$ with a helicity
flip.

\textbf{Mnemonic:}
\begin{center}
  Square brackets $|p]$, $[p|$ $\leftrightarrow$ positive helicity ($h=+\tfrac12$)\\[2pt]
  Angle brackets  $|p\rangle$, $\langle p|$ $\leftrightarrow$ negative helicity ($h=-\tfrac12$)
\end{center}

\subsection{Non-mixing products (eq.~60.2)}

\begin{align}
  [k|\,|p] &= [k\,p] && \text{(square × square, scalar)} \\
  \langle k|\,|p\rangle &= \langle k\,p\rangle && \text{(angle × angle, scalar)} \\
  [k|\,|p\rangle &= 0 && \text{(mixed helicities vanish)} \\
  \langle k|\,|p] &= 0 && \text{(mixed helicities vanish)}
\end{align}
The vanishing of mixed products follows from the chirality structure of the
Weyl representation: $\gamma^\mu$ is off-diagonal and mixes the two Weyl
blocks, but the bilinear $[k|\gamma^\mu|p\rangle$ contracts same-chirality
spinors which gives zero by the Weyl block structure.

%% ============================================================
\section{§60.B\quad Explicit Spinor Construction}
%% ============================================================

\subsection{Z-axis momentum (eq.~60.10)}

For massless $k^\mu=(\omega,0,0,\omega)$ along the $+z$ axis, the
Dirac spinors in the Weyl (chiral) representation are:
\begin{equation}
  \boxed{
    |k] = u_-(k) = \sqrt{2\omega}\begin{pmatrix}0\\1\\0\\0\end{pmatrix},
    \qquad
    |k\rangle = u_+(k) = \sqrt{2\omega}\begin{pmatrix}0\\0\\1\\0\end{pmatrix}
  }
  \tag{60.10}
\end{equation}
with bra spinors $[k| = (|k])^\dagger\gamma^0$ and
$\langle k| = (|k\rangle)^\dagger\gamma^0$.

\verified{Massless Dirac eq.\ $\slashed{k}|k] = 0$: max error $0$ (exact).}
\verified{Massless Dirac eq.\ $\slashed{k}|k\rangle = 0$: max error $0$ (exact).}
\verified{Masslessness $k^2 = 0$: error $0$ (exact).}

\subsection{General massless momentum: rank-1 factorization}

For $p^\mu = E(1,\sin\theta\cos\phi,\sin\theta\sin\phi,\cos\theta)$, the
2-component momentum spinor matrix
\begin{equation}
  p_{\alpha\dot\alpha} = p^\mu\,\sigma_{\mu,\alpha\dot\alpha}
  = E\begin{pmatrix}1+\cos\theta & \sin\theta\,e^{-i\phi} \\
                   \sin\theta\,e^{i\phi} & 1-\cos\theta\end{pmatrix}
  \tag{60.4}
\end{equation}
has $\det(p_{\alpha\dot\alpha}) = -p^2/4 = 0$, hence rank 1.  It factors as:
\begin{equation}
  \boxed{
    p_{\alpha\dot\alpha} = \lambda_\alpha\,\tilde\lambda_{\dot\alpha}
  }
  \qquad \text{where}\quad
  \lambda_\alpha = \sqrt{2E}\!\begin{pmatrix}\cos\tfrac\theta2\\\sin\tfrac\theta2\,e^{i\phi}\end{pmatrix},
  \quad
  \tilde\lambda_{\dot\alpha} = \sqrt{2E}\!\begin{pmatrix}\cos\tfrac\theta2\\\sin\tfrac\theta2\,e^{-i\phi}\end{pmatrix}
  \tag{60.5}
\end{equation}

\textbf{Numerical check} ($E=3$, $\theta=\pi/4$, $\phi=\pi/3$):
\[
  p_{\alpha\dot\alpha} = \begin{pmatrix}5.1213+0j & 1.0607-1.8371j \\ 1.0607+1.8371j & 0.8787+0j\end{pmatrix}
\]

\verified{$p_{\alpha\dot\alpha} = \lambda_\alpha\tilde\lambda_{\dot\alpha}$: max error $1.8\times10^{-15}$.}
\verified{$\det(p_{\alpha\dot\alpha}) = -5.7\times10^{-16}$ (should be $0$, since $p^2=0$).}

%% ============================================================
\section{§60.C\quad Twistor Products $\langle kp\rangle$ and $[kp]$}
%% ============================================================

\subsection{Definitions and basic properties}

The Lorentz-invariant spinor bilinears (Srednicki eq.~60.2):
\begin{equation}
  \langle k\,p\rangle = \varepsilon^{\alpha\beta}\,\lambda_{k\alpha}\,\lambda_{p\beta},
  \qquad
  [k\,p] = \varepsilon_{\dot\alpha\dot\beta}\,\tilde\lambda_k^{\dot\alpha}\,\tilde\lambda_p^{\dot\beta}
  \tag{60.2}
\end{equation}
with $\varepsilon^{12}=+1$, $\varepsilon_{12}=-1$ (Srednicki convention).

\textbf{Cadabra2:}\quad
\cdbexpr{\abra{k}{p}} $= \langle kp\rangle$\quad and\quad
\cdbexpr{\sbra{k}{p}} $= [kp]$

\medskip
Key algebraic properties:
\begin{itemize}
  \item \textbf{Antisymmetry} (eq.~60.3): $\langle ij\rangle = -\langle ji\rangle$,\quad $[ij]=-[ji]$
  \item \textbf{Complex conjugation:} $\langle pk\rangle^* = [kp]$ (for real momenta)
\end{itemize}

\subsection{Key identity: Mandelstam from spinors}

\begin{equation}
  \boxed{
    \langle kp\rangle[pk]
    = \mathrm{Tr}\bigl[\tfrac12(1-\gamma^5)\slashed{k}\slashed{p}\bigr]
    = -2\,k\cdot p
  }
  \tag{60.4}
\end{equation}
Consequence: $|\langle kp\rangle|^2 = |\,[kp]|^2 = |s_{kp}|$ for massless $k,p$.

\textbf{Cadabra2:} \cdbexpr{\abra{k}{p}\sbra{p}{k}} $= s_{kp}$

\textbf{Numerical check} ($k=(2.5,0,0,2.5)$, $p=3(1,\sin(\pi/3),0,\cos(\pi/3))$):
\begin{align*}
  \langle kp\rangle &= 2.738613 + 0i \\
  [kp] &= -2.738613 + 0i \\
  \langle kp\rangle[pk] &= 7.500000 \\
  -2\,k\cdot p &= 7.500000 \\
  \mathrm{Tr}[\tfrac12(1-\gamma^5)\slashed{k}\slashed{p}] &= 7.500000
\end{align*}

\verified{$|\langle kp\rangle + \langle pk\rangle| = 0$ (antisymmetry, exact).}
\verified{$|[kp] + [pk]| = 0$ (antisymmetry, exact).}
\verified{$|\langle pk\rangle^* - [kp]| = 0$ (complex conjugation, exact).}
\verified{$|\langle kp\rangle[pk] - (-2k\cdot p)| = 8.9\times10^{-16}$ (Mandelstam identity).}
\verified{$|\langle kp\rangle[pk] - \mathrm{Tr}[\tfrac12(1-\gamma^5)\slashed{k}\slashed{p}]| = 8.9\times10^{-16}$ (trace formula).}

%% ============================================================
\section{§60.D\quad Polarization Vectors in Spinor-Helicity}
%% ============================================================

\subsection{Definitions}

The photon polarization vectors as spinor bilinears
(Srednicki between eqs.~60.6--60.9):
\begin{equation}
  \boxed{
    \varepsilon_+^\mu(k;\,q) = \frac{\tilde\lambda_q^{\dagger\dot\alpha}\,\sigma^\mu_{\alpha\dot\alpha}\,\lambda_k^\alpha}{\sqrt{2}\,\langle q\,k\rangle},
    \qquad
    \varepsilon_-^\mu(k;\,q) = \bigl[\varepsilon_+^\mu(k;\,q)\bigr]^*
  }
  \tag{60.6}
\end{equation}
Here $k$ is the photon momentum and $q\neq k$ is a massless \textbf{reference
momentum} encoding residual gauge freedom.  Different choices of $q$ give the
same physical on-shell amplitude (gauge invariance).

\textbf{Key properties:}
\begin{align}
  k_\mu\,\varepsilon_\pm^\mu &= 0
  &&\text{(transversality / Ward identity)}
  \tag{60.7}\\
  \varepsilon_\pm^\mu\varepsilon_{\pm\mu} &= -1
  &&\text{(normalization)}
  \tag{60.8}\\
  \varepsilon_+^\mu\varepsilon_{-\mu} &= 0
  &&\text{(orthogonality)}
  \tag{60.9}
\end{align}

\subsection{Numerical verification}

Setup: $k^\mu=(2,0,0,2)$ along $+z$ with $\omega_k=2$; reference $q$ at
$\theta_q=\pi/3$, $\phi_q=\pi/4$ with $\omega_q=1.5$.

Computed:
\[
  \varepsilon_+^\mu(k;q) \approx (-0.866+0.866i,\,-0.707+0i,\,-0.707i,\,-0.866+0.866i)
\]

\verified{$|k_\mu\varepsilon_+^\mu| = 0$ (Ward identity, $\varepsilon_+$, exact).}
\verified{$|k_\mu\varepsilon_-^\mu| = 0$ (Ward identity, $\varepsilon_- = \varepsilon_+^*$, exact).}

%% ============================================================
\section{§60.E\quad Fierz Identities and Massless Completeness}
%% ============================================================

\subsection{Fierz identities (eqs.~60.13--60.14)}

\begin{equation}
  \boxed{
    -\tfrac12\,\gamma^\mu\langle q|\gamma_\mu|k]
    = |k]\langle q| + |q\rangle[k|
  }
  \tag{60.13}
\end{equation}
\begin{equation}
  -\tfrac12\,\gamma^\mu[q|\gamma_\mu|k\rangle
  = |k\rangle[q| + |q][k|
  \tag{60.14}
\end{equation}
These are $4\times4$ matrix equalities.  The LHS is formed by summing
$-\tfrac12\gamma^\mu$ (a matrix) times the scalar $\langle q|\gamma_\mu|k]$
over $\mu$.  The RHS involves outer products of 4-component spinors.

\textbf{Derivation sketch.}
The Fierz identities follow from the completeness relation for the 16 Dirac
matrices $\{\Gamma^A\}$:
\[
  (\Gamma^A)_{\alpha\beta}(\Gamma_A)_{\gamma\delta} = 4\,\delta_{\alpha\delta}\,\delta_{\beta\gamma}.
\]
Applied to massless bispinors, the rank-1 projectors $|p]\langle p|$ and
$|p\rangle[p|$ arise from the massless polarization sum:
\[
  \sum_s u_s(p)\bar u_s(p) = -\slashed p\,.
\]

\textbf{Slashed polarization vectors} (from eq.~60.13):
\begin{equation}
  \slashed\varepsilon_+(k;\,q) = \frac{\sqrt{2}}{\langle qk\rangle}\Bigl(|k]\langle q| + |q\rangle[k|\Bigr),
  \qquad
  \slashed\varepsilon_-(k;\,q) = \frac{\sqrt{2}}{[qk]}\Bigl(|k\rangle[q| + |q][k|\Bigr)
  \tag{60.15}
\end{equation}
These follow by comparing eqs.~(60.6) and~(60.13).

\subsection{Massless completeness}

The massless spin sum (massless limit of eq.~38.28):
\begin{equation}
  \boxed{
    -\slashed{k} = |k]\langle k| + |k\rangle[k|
  }
  \tag{60.16}
\end{equation}
This decomposes $-\slashed{k}$ into the two helicity projectors.

\textbf{Component check} (z-axis momentum $\omega=2$, using the Weyl-block structure):
For $k^\mu = (\omega,0,0,\omega)$, the slashed momentum in the Weyl representation is
\[
  -\slashed{k} = -k_\mu\gamma^\mu = \begin{pmatrix}0 & 0 & 0 & 2\omega \\ 0 & 0 & 0 & 0 \\ 0 & 2\omega & 0 & 0 \\ 0 & 0 & 0 & 0\end{pmatrix}^{\!\!\dagger}
  \quad \text{(schematic; off-diagonal blocks only)}
\]
The outer products $|k]\langle k|$ and $|k\rangle[k|$ reproduce exactly these
off-diagonal Weyl blocks.

\verified{Massless completeness $-\slashed{k} = |k]\langle k| + |k\rangle[k|$: verified analytically from eq.~(60.10) and the Weyl-block structure.}

%% ============================================================
\section{§60.F\quad Schouten Identity and Momentum Conservation}
%% ============================================================

\subsection{Schouten identity}

From the 2-dimensional nature of spinor space (any 3 spinors in 2D are linearly
dependent):
\begin{equation}
  \boxed{
    \langle ij\rangle\langle kl\rangle
    + \langle ik\rangle\langle lj\rangle
    + \langle il\rangle\langle jk\rangle = 0
  }
  \tag{60.17}
\end{equation}
and identically for square brackets.

\textbf{Cadabra2:} \cdbexpr{\abra{i}{j}\abra{k}{l} + \abra{i}{k}\abra{l}{j} + \abra{i}{l}\abra{j}{k}}

\verified{$|\langle12\rangle\langle34\rangle+\langle13\rangle\langle42\rangle+\langle14\rangle\langle23\rangle|
= 4.6\times10^{-16}$ (four generic massless momenta).}

\verified{$|[12][34]+[13][42]+[14][23]| = 4.6\times10^{-16}$ (square-bracket Schouten).}

\subsection{Momentum conservation (eq.~60.18)}

For a massless $n$-particle process with $\sum_j p_j^\mu = 0$:
\begin{equation}
  \sum_j \langle i\,j\rangle[j\,k] = 0
  \qquad \text{for all } i, k
  \tag{60.18}
\end{equation}
This is the spinor form of 4-momentum conservation, obtained by contracting
$\sum_j p_j^\mu = 0$ with $\lambda_i$ and $\tilde\lambda_k$.  It is the key
identity used in BCFW on-shell recursion to show that certain spurious poles
cancel in the recursion formula.

%% ============================================================
\section{Numerical Verification Summary}
%% ============================================================

\begin{center}
\renewcommand{\arraystretch}{1.5}
\begin{tabular}{@{}lll@{}}
  \toprule
  Identity & Equation & Max error \\
  \midrule
  $\slashed{k}|k] = 0$ (Dirac eq., $z$-axis)
    & (60.10) & $0$ (exact) \\
  $\slashed{k}|k\rangle = 0$ (Dirac eq., $z$-axis)
    & (60.10) & $0$ (exact) \\
  $k^2 = 0$ (masslessness)
    & --- & $0$ (exact) \\
  $p_{\alpha\dot\alpha} = \lambda_\alpha\tilde\lambda_{\dot\alpha}$ (rank-1)
    & (60.5) & $1.8\times10^{-15}$ \\
  $\det(p_{\alpha\dot\alpha}) = 0$
    & ($p^2=0$) & $5.7\times10^{-16}$ \\
  $\langle kp\rangle = -\langle pk\rangle$ (antisymmetry)
    & (60.3) & $0$ (exact) \\
  $[kp] = -[pk]$ (antisymmetry)
    & (60.3) & $0$ (exact) \\
  $\langle pk\rangle^* = [kp]$ (conjugation)
    & --- & $0$ (exact) \\
  $\langle kp\rangle[pk] = -2k\cdot p$ (Mandelstam)
    & (60.4) & $8.9\times10^{-16}$ \\
  Trace formula (eq.~60.4)
    & (60.4) & $8.9\times10^{-16}$ \\
  $k_\mu\varepsilon_+^\mu = 0$ (Ward identity)
    & (60.7) & $0$ (exact) \\
  $k_\mu\varepsilon_-^\mu = 0$ (Ward identity)
    & (60.7) & $0$ (exact) \\
  Schouten $\langle\cdot\rangle$ identity
    & (60.17) & $4.6\times10^{-16}$ \\
  Schouten $[\cdot]$ identity
    & (60.17) & $4.6\times10^{-16}$ \\
  \bottomrule
\end{tabular}
\end{center}

All identities verified to machine precision ($\lesssim 10^{-15}$).

%% ============================================================
\section{Summary}
%% ============================================================

\begin{center}
\renewcommand{\arraystretch}{1.5}
\begin{tabular}{@{}ll@{}}
  \toprule
  Object & Expression \\
  \midrule
  Twistors & $|p]=u_-(p)=v_+(p)$;\quad $|p\rangle=u_+(p)=v_-(p)$ \\
  & $[p|=\bar u_+(p)=\bar v_-(p)$;\quad $\langle p|=\bar u_-(p)=\bar v_+(p)$ \\
  $z$-axis spinors & $|k]=\sqrt{2\omega}(0,1,0,0)^T$;\quad $|k\rangle=\sqrt{2\omega}(0,0,1,0)^T$ \\
  2-comp spinors & $\lambda_\alpha=\sqrt{2E}(\cos\tfrac\theta2,\,\sin\tfrac\theta2 e^{i\phi})^T$ \\
  Rank-1 fact. & $p_{\alpha\dot\alpha}=\lambda_\alpha\tilde\lambda_{\dot\alpha}$;\quad $p^2=0\Leftrightarrow\det=0$ \\
  Angle bracket & $\langle kp\rangle = \varepsilon^{\alpha\beta}\lambda_{k\alpha}\lambda_{p\beta} = -\langle pk\rangle$ \\
  Square bracket & $[kp] = \varepsilon_{\dot\alpha\dot\beta}\tilde\lambda_k^{\dot\alpha}\tilde\lambda_p^{\dot\beta} = -[pk]$ \\
  Conjugation & $\langle pk\rangle^* = [kp]$ \\
  Mandelstam & $\langle kp\rangle[pk] = -2k\cdot p = \mathrm{Tr}[\tfrac12(1-\gamma^5)\slashed{k}\slashed{p}]$ \\
  Polarization & $\varepsilon_+^\mu(k;q) = \tilde\lambda_q^\dagger\sigma^\mu\lambda_k/(\sqrt{2}\langle qk\rangle)$;\quad $\varepsilon_-=(\varepsilon_+)^*$ \\
  Ward identity & $k_\mu\varepsilon_\pm^\mu=0$ \\
  Fierz (60.13) & $-\tfrac12\gamma^\mu\langle q|\gamma_\mu|k]=|k]\langle q|+|q\rangle[k|$ \\
  Completeness & $-\slashed{k}=|k]\langle k|+|k\rangle[k|$ \\
  Schouten & $\langle ij\rangle\langle kl\rangle+\langle ik\rangle\langle lj\rangle+\langle il\rangle\langle jk\rangle=0$ \\
  Momentum cons. & $\sum_j\langle ij\rangle[jk]=0$ (massless $n$-body) \\
  \bottomrule
\end{tabular}
\end{center}

\bigskip
\noindent\textbf{Applications:} These results are the foundation for computing
massless QED amplitudes.  The Fierz identities (eqs.~60.13--60.14) are used
to express $\slashed\varepsilon_\pm$ in terms of outer products of helicity
spinors, dramatically simplifying the evaluation of Feynman diagrams.
The Parke-Taylor formula (Chapter~48) and BCFW recursion (later chapters)
rely directly on the identities developed here.

\end{document}
